\documentclass[aspectratio=169]{beamer}
\usetheme{Madrid}
\usecolortheme{beaver}
\usepackage[utf8]{inputenc}
\usepackage[T1]{fontenc}
\usepackage{tikz}
\usepackage{amssymb}
\usepackage{graphicx}
\setbeamertemplate{navigation symbols}{}
\setbeamercolor{title}{fg=white,bg=cyan!60!black}
\setbeamercolor{frametitle}{fg=white,bg=cyan!50!black}
\setbeamercolor{block title}{bg=cyan!15,fg=black}
\setbeamercolor{block body}{bg=cyan!5,fg=black}

\title{G9 — Strings \& Dialogue}
\subtitle{Player names, dialogue trees, high-score tables}
\author{ByteForge --- Grade 7 --- Game Maker --- Week Overview}
\date{}

\begin{document}

\begin{frame}
    \titlepage
\end{frame}

\begin{frame}{This Week's Plan}
    \begin{itemize}
        \item Lesson 1: Player Name \& ASCII
        \item Lesson 2: Dialogue Tree
        \item Lesson 3: Shout Mode
        \item[$\bigstar$] BOSS --- RPG Dialogue Scene (\texttt{lessonBoss.tex})
    \end{itemize}
\end{frame}

\begin{frame}{Real-World Example}
    \Large
    Any RPG's dialogue choices ("Attack" / "Defend" / "Run") branching to different outcomes, and a game's online leaderboard listing name+score pairs, are both exactly this chapter's ideas in shipped games. Minecraft villager trading menus are a real dialogue tree too — compare your held item's name against what the villager wants, branch to "trade accepted" or not.
\end{frame}


\begin{frame}{Let Them Discover}
    \Large
    Play (or recall) a game with dialogue choices (or a Minecraft villager trade). Pick one branch and trace it: what text comparison do you think the code is making to decide which reply you see?
\end{frame}

\begin{frame}{Files in This Folder}
    \small
    \begin{itemize}
        \item \texttt{lesson1.tex} --- Player Name \& ASCII
        \item \texttt{lesson2.tex} --- Dialogue Tree
        \item \texttt{lesson3.tex} --- Shout Mode
        \item \texttt{lessonBoss.tex} --- BOSS: RPG Dialogue Scene
    \end{itemize}
    \vspace{0.5em}
    \normalsize Each lesson file is its own short deck --- open the one for
    the specific day you're teaching, not this whole overview.
\end{frame}

\end{document}
